\chapter*{Résumé}
% Ajouter manuellement cette section à la table des matières
\addcontentsline{toc}{chapter}{Résumé}

La directive \textit{Solvabilité II}\ impose aux organismes d'assurance une évaluation prospective de leurs engagements financiers, nécessitant le déploiement de modèles de gestion actif-passif stochastiques. Ces projections requièrent une puissance de calcul considérable, rendant le traitement unitaire des contrats d'assurance vie opérationnellement prohibitif. Pour atténuer cette complexité algorithmique, la réduction de la dimension du portefeuille par un processus d'agrégation en \textit{Model Points} constitue une solution adaptée. Toutefois, un regroupement excessif risque de lisser les garanties financières, altérant ainsi l'évaluation du \textit{Best Estimate} (BE) et la restitution de la structure du bilan face aux chocs de marché, faussant potentiellement le Capital de Solvabilité Requis (SCR).

L'objectif de ce mémoire est d'évaluer l'influence de différentes méthodes d'agrégation de données sur la détermination du BE et du SCR. La démarche repose sur la construction d'un portefeuille de référence de 1\,000\,000 de polices, représentatif d'un assureur vie moyen sur le marché français.

Ce portefeuille permet de confronter plusieurs techniques de regroupement : le découpage par classes (\textit{Banding}), les algorithmes de \textit{Clustering} (\textit{K-Means} et \textit{HDBSCAN}) ainsi qu'une approche par appariement de flux (\textit{Cash-Flow Matching}). L'étude compare ces méthodes afin d'identifier celle permettant de minimiser l'erreur résiduelle sur les indicateurs de référence tout en optimisant les temps de calcul.

L'analyse tend à démontrer que l'allègement de la dimensionnalité permet de diviser les temps de traitement par un facteur cinquante. Face à des chocs de taux d'intérêt asymétriques, les algorithmes spatiaux, et particulièrement la méthode \textit{K-Means}, offrent le meilleur compromis opérationnel. En préservant la granularité des taux minimums garantis et de l'ancienneté des contrats, ces modèles capturent avec précision la mécanique financière du coût d'opportunité inhérent au risque de rachat. À l'inverse, l'approche fondée sur le \textit{Cash-Flow Matching} tend à lisser l'information fiscale, conduisant à une surévaluation mécanique du SCR estimé.

\vspace{0.5cm}
\noindent \textbf{Mots-clés :} assurance vie, ALM, Solvabilité II, agrégation, \textit{Model Points}, \textit{Machine Learning}, \textit{K-Means}, analyse de sensibilités, SCR, \textit{Best Estimate}.